% Gerado por regressao/06_tabelas_latex.R (projeto Modelagem). Não editar à mão.
\begin{sidewaystable}
\caption{Robustez R13 — modelo B com erro-padrão de Driscoll-Kraay}\label{tab:robustez_R13}
\centering
\footnotesize
\setlength{\tabcolsep}{3pt}
\begin{tabular}{@{}>{\raggedright\arraybackslash}p{\dimexpr\linewidth-594pt\relax}>{\centering\arraybackslash}p{60pt}>{\centering\arraybackslash}p{60pt}>{\centering\arraybackslash}p{60pt}>{\centering\arraybackslash}p{60pt}>{\centering\arraybackslash}p{60pt}>{\centering\arraybackslash}p{60pt}>{\centering\arraybackslash}p{60pt}>{\centering\arraybackslash}p{60pt}>{\centering\arraybackslash}p{60pt}@{}}
\hline
\textbf{Variável} & \textbf{\hspace{0pt}Despesa total} & \textbf{\hspace{0pt}Saúde e Saneamento} & \textbf{\hspace{0pt}Educação e Cultura} & \textbf{\hspace{0pt}Habitação e Urbanismo} & \textbf{\hspace{0pt}Assistência Social} & \textbf{\hspace{0pt}Transporte} & \textbf{\hspace{0pt}Administração} & \textbf{\hspace{0pt}Agricultura} & \textbf{\hspace{0pt}Comunicações} \\ \hline
\multicolumn{10}{@{}l}{\textit{Modelo B (coalizões)}} \\
Prefeito na coalizão & 0,0037\textsuperscript{***} & 0,0009 & 0,0033 & 0,0160\textsuperscript{***} & 0,0009 & 0,0114 & 0,0070\textsuperscript{***} & 0,0110\textsuperscript{*} & $-$0,0022 \\
do governador & (0,0014) & (0,0023) & (0,0022) & (0,0051) & (0,0031) & (0,0094) & (0,0019) & (0,0064) & (0,0085) \\
Prefeito na coalizão & $-$0,0010 & 0,0001 & $-$0,0025 & $-$0,0163\textsuperscript{***} & $-$0,0058\textsuperscript{*} & 0,0418\textsuperscript{***} & 0,0043 & 0,0045 & 0,0075 \\
do presidente & (0,0022) & (0,0020) & (0,0041) & (0,0058) & (0,0032) & (0,0101) & (0,0033) & (0,0090) & (0,0271) \\
\hline
Observações & 70.669 & 70.526 & 70.550 & 69.574 & 70.484 & 54.422 & 70.553 & 64.464 & 12.360 \\
Municípios & 5.568 & 5.568 & 5.568 & 5.564 & 5.568 & 5.205 & 5.568 & 5.436 & 1.896 \\
R\textsuperscript{2} & 0,390 & 0,163 & 0,196 & 0,036 & 0,069 & 0,019 & 0,075 & 0,017 & 0,013 \\
\hline
\end{tabular}
\fonte{Elaboração própria.}
\nota{Despesa e receitas \textit{per capita}, em R\$ de 2025, em logaritmo (\autoref{sec:especificacao}). O coeficiente de uma \textit{dummy} se lê como variação percentual aproximada (exata: 100 $\times$ (exp(b) $-$ 1)). Só modelo B: efeitos fixos de município e de ano. Erros-padrão de Driscoll-Kraay entre parênteses. \textsuperscript{***}p$<$0,01; \textsuperscript{**}p$<$0,05; \textsuperscript{*}p$<$0,1.}
\end{sidewaystable}
